%+TITLE: Representaciones de SO(3)
%+DESCRIPTION: Discutimos las representaciones del grupo de rotaciones antes de discutir repreaentaciones en general.
\documentclass[11pt,a4paper]{article}

\usepackage[utf8]{inputenc}
\usepackage[T1]{fontenc}
\usepackage[spanish,es-nodecimaldot]{babel}
\usepackage{lmodern}
\usepackage{amsmath,amssymb,mathtools}
\usepackage{graphicx}
\usepackage{xcolor}
\usepackage{geometry}
\usepackage{enumitem}
\usepackage{microtype}

\geometry{margin=2.3cm}
\setlength{\parindent}{0pt}
\setlength{\parskip}{0.65em}
\setlist[itemize]{leftmargin=1.6em,itemsep=0.25em,topsep=0.25em}

\definecolor{title-red}{RGB}{165,25,20}
\definecolor{concept-green}{RGB}{20,125,45}
\newcommand{\concept}[1]{\textcolor{concept-green}{#1}}

\newcommand{\ket}[1]{\left|#1\right\rangle}
\newcommand{\bra}[1]{\left\langle#1\right|}
\newcommand{\braket}[2]{\left\langle #1\middle|#2\right\rangle}
\newcommand{\hc}{\mathrm{h.c.}}
\newcommand{\Tr}{\operatorname{Tr}}
\newcommand{\Diffs}{\mathrm{Diffs}}

\begin{document}

Ahora buscamos las representaciones finito-dimensionales del grupo de rotaciones. En realidad lo que haremos es buscar representaciones del álgebra de Lie y tomaremos su exponencial.

Una representación del álgebra de Lie es un conjunto de operadores lineales sobre el espacio $V$ tal que su conmutador satisface el álgebra.

Una representación del álgebra de Lie me permite encontrar una representación del grupo localmente.

\[
  g_1=e^{i\theta_1^aG^a},
  \qquad
  g_2=e^{i\theta_2^aG^a},
\]
\[
  g_1g_2
  =e^{i\left(
      (\theta_1^a+\theta_2^a)G^a
      +\frac{i}{2}\theta_1^a\theta_2^b[G^a,G^b]
      +\cdots
    \right)}.
\]
De aquí se puede ver que
\[
  R(g_1)R(g_2)=R(g_1g_2),
  \qquad
  R(g)=e^{i\theta^aR(G^a)}.
\]

Para encontrar las representaciones de dimensión finita del álgebra, usamos un hecho del álgebra lineal:

Dos matrices son mutuamente diagonalizables si y sólo si conmutan.

En el caso de las rotaciones, ninguna pareja de operadores $J_1,J_2,J_3$ conmuta. Sólo uno de ellos puede ser diagonal. Sin perder generalidad, escojamos $J_3$. Que sea diagonal quiere decir que los elementos de la base son autovectores de $J_3$:
\[
  J_3|m\rangle=m|m\rangle.
\]

Para ver cómo actúan los otros generadores, definimos
\[
  J_\pm=(J_1\pm iJ_2).
\]
Entonces
\begin{equation}
  [J_3,J_+]
  =[J_3,J_1+iJ_2]
   =(iJ_2+J_1)
   =J_+.
\end{equation}
Análogamente,
\[
  [J_3,J_-]=-J_-.
\]

De aquí tenemos que
\begin{align*}
  J_3J_+|m\rangle
  &=([J_3,J_+]+J_+J_3)|m\rangle\\
  &=J_+|m\rangle+J_+J_3|m\rangle\\
  &=(m+1)J_+|m\rangle.
\end{align*}
Análogamente,
\[
  J_3J_-|m\rangle=(m-1)J_-|m\rangle.
\]

Esto quiere decir que $J_+|m\rangle$ es un autovector de $J_3$ con autovalor $m+1$, y $J_-|m\rangle$ es un autovector de $J_3$ con autovalor $m-1$.

Si la representación es finita, $m$ debe tener algún rango finito. Llamemos $j$ al valor máximo que puede tomar $m$, entonces
\[
  J_+|j\rangle=0.
\]

El estado $|j-1\rangle$ debe estar normalizado, escribimos
\[
  J_-|j\rangle=\alpha_0|j-1\rangle.
\]
Entonces
\[
  |\alpha_0|^2\langle j-1|j-1\rangle
  =\langle j|J_+J_-|j\rangle
  =\langle j|([J_+,J_-]+J_-J_+)|j\rangle
\]

Pero
\[
  [J_+,J_-]=[J_1+iJ_2,J_1-iJ_2]=2J_3,
\]
tal que
\[
  |\alpha_0|^2=2\langle j|J_3|j\rangle=2j
  \quad\Longrightarrow\quad
  \alpha_0=\sqrt{2j}.
\]

Queremos demostrar en general
\[
  J_-|j-n\rangle
  =\sqrt{(2j-n)(n+1)}\,|j-n-1\rangle.
\]

Esto se deduce a partir de
\[
  J_+J_-|j-n\rangle=(2j-n)(n+1)|j-n\rangle,
\]
que demostramos por recurrencia:

\begin{itemize}
  \item Ya demostramos el caso base $n=0$.
  \item Asumimos que se cumple para $n-1$ y lo extendemos a $n$.
\end{itemize}

\begin{align*}
  J_+J_-|j-n\rangle
  &=(2J_3+J_-J_+)|j-n\rangle \\
  &=2(j-n)|j-n\rangle
    +(2j-n+1)^{-1/2}n^{-1/2}
      J_-J_+J_-|j-n+1\rangle \\
  &=2(j-n)|j-n\rangle
    +(2j-n+1)^{-1/2}n^{-1/2}
      (2j-n+1)n\,J_-|j-n+1\rangle \\
  &=2(j-n)|j-n\rangle+(2j-n+1)n|j-n\rangle \\
  &=(2j-n+2jn-n^2)|j-n\rangle \\
  &=(2j-n)(n+1)|j-n\rangle.
\end{align*}

Entonces
\begin{align*}
  |\alpha_n|^2
  &=|\alpha_n|^2\langle j-n-1|j-n-1\rangle \\
  &=\langle j-n|J_+J_-|j-n\rangle \\
  &=(2j-n)(n+1).
\end{align*}

Escribiendo en términos de $m$,
\[
  J_-|m\rangle
  =\sqrt{(j+m)(j-m+1)}\,|m-1\rangle.
\]
Análogamente,
\[
  J_+|m\rangle
  =\sqrt{(j-m)(j+m+1)}\,|m+1\rangle.
\]

Si $j$ es cualquier número real, la representación no es necesariamente finita:
\[
  |1.2\rangle\xrightarrow{J_-}|0.2\rangle
  \xrightarrow{J_-}|-0.8\rangle
  \xrightarrow{J_-}|-1.8\rangle
  \xrightarrow{J_-}\cdots
\]

Nos da infinitos vectores. Pero si $j$ toma algunos valores especiales, la normalización $\sqrt{(2j-n)(n+1)}$ se hace cero luego de un número finito de ``bajadas'':
\[
  n=2j.
\]

Como $n$ es entero, vemos que esto se puede cumplir sólo si $2j$ es entero.

El número de vectores generados es $2j+1$.

Para $j=0$ hay un solo vector y corresponde a la representación trivial.

Para $j=1/2$ hay dos vectores:
\[
  J_3|1/2\rangle=\frac12|1/2\rangle,
  \qquad
  J_3|-1/2\rangle=-\frac12|-1/2\rangle,
\]
\[
  J_+|-1/2\rangle=|1/2\rangle,
  \qquad
  J_-|1/2\rangle=|-1/2\rangle.
\]

Estos forman una base. Cualquier vector en $V$ se escribe
\[
  \alpha|1/2\rangle+\beta|-1/2\rangle
  \quad\text{o}\quad
  \begin{pmatrix}\alpha\\ \beta\end{pmatrix}.
\]

Este vector se llama \concept{espinor de Pauli}. Por ejemplo, los electrones transforman en esta representación bajo rotaciones.

\[
  J_3=\frac12
  \begin{pmatrix}
    1&0\\
    0&-1
  \end{pmatrix},
  \qquad
  J_+=
  \begin{pmatrix}
    0&1\\
    0&0
  \end{pmatrix},
  \qquad
  J_-=
  \begin{pmatrix}
    0&0\\
    1&0
  \end{pmatrix}.
\]
Por tanto,
\[
  J_1=\frac12(J_++J_-)
  =\frac12
  \begin{pmatrix}
    0&1\\
    1&0
  \end{pmatrix},
  \qquad
  J_2=\frac{1}{2i}(J_+-J_-)
  =\frac12
  \begin{pmatrix}
    0&-i\\
    i&0
  \end{pmatrix}.
\]

\end{document}
